\documentclass[conference]{IEEEtran}
\IEEEoverridecommandlockouts
\usepackage{cite}
\usepackage{amsmath,amssymb,amsfonts,amsthm}
\usepackage{algorithmic}
\usepackage{graphicx}
\usepackage{textcomp}
\usepackage{xcolor}
\usepackage{booktabs}
\usepackage{multirow}

\newtheorem{theorem}{Theorem}
\newtheorem{lemma}{Lemma}
\newtheorem{remark}{Remark}

\def\BibTeX{{\rm B\kern-.05em{\sc i\kern-.025em b}\kern-.08em
    T\kern-.1667em\lower.7ex\hbox{E}\kern-.125emX}}

\begin{document}

\title{Structural Reasoning versus Statistical Estimation for Learned Matrix Chain Optimization\\
{\footnotesize \textsuperscript{}}
}

\author{
\IEEEauthorblockN{1\textsuperscript{st} Ch. Abhilash}
\IEEEauthorblockA{\textit{Department of CSE} \\
\textit{VNR VJIET}\\
Hyderabad, India \\
abhilash00042@gmail.com}
\and
\IEEEauthorblockN{2\textsuperscript{nd} Dakshatha}
\IEEEauthorblockA{\textit{Department of CSE} \\
\textit{VNR VJIET}\\
Hyderabad, India \\
dklakshatha18@gmail.com}
\and
\IEEEauthorblockN{3\textsuperscript{rd} Shaun Angel}
\IEEEauthorblockA{\textit{Department of CSE} \\
\textit{VNR VJIET}\\
Hyderabad, India \\
dshaunangel@gmail.com}
\and
\IEEEauthorblockN{4\textsuperscript{th} Abhinav}
\IEEEauthorblockA{\textit{Department of CSE} \\
\textit{VNR VJIET}\\
Hyderabad, India \\
abhinavkanna2627@gmail.com}
\and
\IEEEauthorblockN{5\textsuperscript{th} Maneeshwar}
\IEEEauthorblockA{\textit{Department of CSE} \\
\textit{VNR VJIET}\\
Hyderabad, India \\
maneeshwarkarne@gmail.com}
}

\maketitle

\begin{abstract}
The Matrix Chain Multiplication (MCM) problem requires finding the most efficient
sequence for multiplying a chain of matrices. Classical dynamic programming (DP)
solves it exactly in $O(n^3)$ time, which becomes a practical bottleneck in systems
that must evaluate many chains repeatedly, such as deep learning compilers and
database query optimizers. This paper studies whether learned models can
approximate the DP solver with negligible error, and---more importantly---%
\emph{how the problem is framed} determines whether learned predictions are even
mathematically meaningful. Two paradigms are compared under identical conditions:
\emph{structural reasoning}, in which a Graph Neural Network (GraphMCMNet) and a
Pointer Network (PointerMCMNet) predict the DP split table and the cost is then
computed exactly from the predicted structure, versus \emph{statistical
estimation}, in which XGBoost and Random Forest ensembles regress the scalar cost
from 145 engineered features. We prove that any split-table prediction induces a
feasible parenthesization, so structural models can never report a cost below the
DP optimum. On a held-out benchmark of 500 synthetic chains ($n{=}5$--$50$) drawn
from four dimension distributions, GraphMCMNet attains 0.029\% MAPE, a
log-scale $R^2$ of 0.999999, and reproduces the exact DP split table in 93\% of
cases, while maintaining 100\% mathematical validity. In contrast, the tree
ensembles---despite extensive feature engineering---predict costs below the DP
lower bound in 78--82\% of cases. Wilcoxon signed-rank tests confirm all pairwise
differences ($p{<}0.001$), with very large effect sizes ($|d|{>}1.0$) between the
structural and statistical families. Trained only on chains of length
$n{\le}35$, GraphMCMNet achieves its \emph{lowest} error (0.019\% MAPE) on
out-of-distribution chains of length $41$--$50$. On our synthetic benchmark,
these results indicate that predicting combinatorial structure, rather than
regressing unconstrained costs, is the decisive design choice for learned MCM
approximation.
\end{abstract}

\begin{IEEEkeywords}
matrix chain multiplication, neural combinatorial optimization, graph neural
network, pointer network, dynamic programming, neural algorithmic reasoning,
structural reasoning
\end{IEEEkeywords}

\section{Introduction}
The Matrix Chain Multiplication (MCM) problem is a cornerstone of dynamic
programming: given $n$ matrices, determine the parenthesization that minimizes
the total number of scalar multiplications. The classical dynamic programming
solution runs in $O(n^3)$ time and $O(n^2)$ space \cite{cormen2009}. Hu and Shing
\cite{hushing1982} later proposed an $O(n \log n)$ algorithm by reducing the
problem to convex polygon partitioning; while asymptotically superior, the
algorithm is considerably more intricate to implement and has seen limited
adoption in practical optimization systems \cite{lopez2023}. Exact optimization
therefore remains computationally non-trivial in large-scale systems where
thousands of chains must be evaluated repeatedly---a routine requirement in deep
learning compilers \cite{tvm2018,tiramisu2021}, large-scale AI training
pipelines, database query optimizers \cite{neo2019,balsa2022}, and tensor
contraction engines.

Prior attempts to accelerate MCM have relied on heuristic search or evolutionary
algorithms \cite{naz2021}, which offer neither optimality guarantees nor
structural feasibility guarantees. A natural question arises: can a neural
network learn to approximate the DP solver with negligible error while
guaranteeing that every prediction corresponds to a \emph{feasible}
parenthesization?

Recent advances in neural combinatorial optimization
\cite{vinyals2015,bello2017}, graph neural networks for structured problems
\cite{cappart2021,bengio2021,vesselinova2020}, and neural algorithmic reasoning
\cite{velickovic2021nar,ibarz2022generalist,velickovic2022clrs} suggest the
answer may be yes---provided the architecture respects the problem's
combinatorial constraints. However, the dominant paradigm in learned cost
modeling \cite{neo2019,balsa2022,tiramisu2021} treats optimization targets as
unconstrained continuous variables, ignoring the discrete feasibility structure
that DP enforces.

This paper makes the following contributions:
\begin{enumerate}
    \item We introduce \textbf{GraphMCMNet}, a graph neural network that maps
    the DP sub-problem lattice to an explicit graph and predicts optimal split
    points, attaining 0.029\% MAPE and 100\% mathematical validity on a
    500-chain held-out benchmark.
    \item We formally prove (Theorem~\ref{thm:validity}) that any model
    predicting per-sub-chain split points induces a feasible parenthesization,
    so its reported cost can never fall below the DP optimum---a guarantee that
    no unconstrained cost regressor can offer.
    \item We provide a controlled head-to-head comparison of structural
    reasoning (GNN and Pointer Network) versus statistical estimation (XGBoost
    and Random Forest with 145 engineered features), observing a
    three-order-of-magnitude accuracy gap on our benchmark, with all
    differences statistically significant ($p{<}0.001$, $|d|{>}1.0$).
    \item We show that the tree-based regressors violate the DP lower bound in
    78--82\% of cases, producing infeasible cost predictions---a failure mode
    that structural architectures eliminate \emph{by construction}.
    \item We demonstrate length extrapolation: trained only on chains with
    $n{\le}35$, GraphMCMNet achieves its lowest error (0.019\% MAPE) on
    out-of-distribution chains 43\% longer than any training instance.
\end{enumerate}

We deliberately restrict our claims to the synthetic benchmark studied here;
Section~\ref{sec:limitations} discusses external validity and the experiments
required to extend these conclusions to production workloads.

The paper is organized as follows. Section~II reviews related work,
Section~III formulates the problem and states the validity guarantee,
Section~IV presents the methodology, Section~V describes the experimental
setup, Section~VI reports results, Section~VII discusses limitations, and
Section~VIII concludes.

\section{Related Work}

\textbf{Classical MCM Algorithms.} Godbole \cite{godbole1973} first formalized
efficient matrix chain computation via the $O(n^3)$ dynamic program also
presented in standard references \cite{cormen2009}. Hu and Shing
\cite{hushing1982} introduced an $O(n \log n)$ algorithm exploiting a reduction
to polygon partitioning. L\'opez et al. \cite{lopez2023} recently revisited the
space of parenthesizations, showing that although all orderings are feasible,
few are ever optimal. Despite its asymptotic advantage, the Hu--Shing algorithm
is rarely deployed in practical optimizers due to implementation complexity; we
discuss its relationship to our approach in Sections~VI-C and~VII.

\textbf{Neural Combinatorial Optimization.} Vinyals et al. \cite{vinyals2015}
introduced Pointer Networks for sequence-to-sequence combinatorial problems,
and Bello et al. \cite{bello2017} extended them with reinforcement learning.
Bahdanau et al. \cite{bahdanau2015} proposed the additive attention mechanism
that underpins pointer decoders, and the Transformer \cite{vaswani2017} further
advanced sequential modeling. These approaches treat problems as flat sequences
and lack explicit structural inductive biases.

\textbf{GNNs for Optimization.} Cappart et al. \cite{cappart2021} surveyed GNNs
in combinatorial optimization, highlighting their ability to encode problem
structure. Bengio et al. \cite{bengio2021} provided a methodological framework
for ML-driven optimization, emphasizing problem-aware encodings. Vesselinova et
al. \cite{vesselinova2020} demonstrated that learning on graphs improves
generalization for combinatorial tasks. Graph Transformers
\cite{rampasek2022gps,muller2023gt} extend attention to arbitrary graph
structure. Our work extends this line by constructing the DP sub-problem graph
as an explicit GNN input.

\textbf{Neural Algorithmic Reasoning.} A growing body of work trains networks
to \emph{execute} classical algorithms rather than regress their outputs.
Veli\v{c}kovi\'c and Blundell \cite{velickovic2021nar} articulated the neural
algorithmic reasoning blueprint; the CLRS-30 benchmark
\cite{velickovic2022clrs} and the generalist algorithmic learner
\cite{ibarz2022generalist} showed that message-passing networks aligned with an
algorithm's intermediate state generalize far better than unaligned models, and
Bounsi et al. \cite{bounsi2024transformers} combined Transformers with
algorithmic reasoners. Our split-table prediction task is an instance of this
paradigm applied to the MCM dynamic program: the sub-problem graph aligns the
network's computation with the DP recurrence, which we conjecture underlies the
strong length extrapolation observed in Section~VI-B.

\textbf{Learned Cost Models.} In query optimization, Neo \cite{neo2019} and
Balsa \cite{balsa2022} train neural models to predict query costs as regression
targets. In compiler optimization, TVM \cite{tvm2018} and Tiramisu
\cite{tiramisu2021} use learned cost models to guide search. A critical
limitation of these approaches is that they regress unconstrained scalars,
offering no guarantee that predictions correspond to feasible solutions. Our
structural approach addresses this gap directly and quantifies the resulting
validity failure empirically.

\section{Problem Formulation and Validity Guarantee}
\label{sec:formulation}

Given a chain of $n$ matrices $A_1, A_2, \ldots, A_n$, where matrix $A_i$ has
dimensions $d_{i-1} \times d_i$, the MCM problem seeks a parenthesization that
minimizes the total number of scalar multiplications. Although the final
product is identical regardless of evaluation order, the cost varies
dramatically: for a 5-matrix chain with dimensions $[10,100,5,50,10,30]$, one
parenthesization requires $10{,}500$ multiplications while another demands over
$200{,}000$. The classical DP recurrence \cite{cormen2009} solves this exactly:
\begin{equation}
m[i][j] =
\begin{cases}
0 & \text{if } i = j \\
\displaystyle\min_{i \le k < j} \bigl\{ m[i][k] + m[k{+}1][j] + d_{i-1} d_k d_j \bigr\} & \text{if } i < j
\end{cases}
\label{eq:dp}
\end{equation}
Here $m[i][j]$ is the minimum cost for sub-chain $A_i \cdots A_j$, and the
optimal split table $s[i][j] = \arg\min_k \{\cdot\}$ records the best split for
each sub-chain. Computing the full solution requires $O(n^3)$ time and $O(n^2)$
space.

\textbf{Central Hypothesis.} Can machine learning models approximate this
solver while preserving solution quality and \emph{feasibility}? We compare two
fundamentally different framings:
\begin{enumerate}
    \item \textbf{Structural Reasoning (Split Prediction):} the model predicts
    a split table $\hat{s}[i][j]$ for every sub-chain; the cost is then
    computed \emph{exactly} from these splits.
    \item \textbf{Statistical Estimation (Cost Regression):} the model directly
    regresses a scalar cost $\hat{c}$ from hand-crafted features, with no
    constraint tying $\hat{c}$ to any parenthesization.
\end{enumerate}

The following result formalizes why the first framing is safe by construction.

\begin{theorem}[Structural Validity Guarantee]
\label{thm:validity}
Let $\hat{s}$ be any predicted split table such that
$\hat{s}[i][j] \in \{i, i{+}1, \ldots, j{-}1\}$ for every sub-chain
$1 \le i < j \le n$. Define the induced cost recursively as
\begin{equation}
\hat{m}[i][j] =
\begin{cases}
0 & i = j\\
\hat{m}[i][k^\ast] + \hat{m}[k^\ast{+}1][j] + d_{i-1} d_{k^\ast} d_j & i < j
\end{cases}
\end{equation}
with $k^\ast = \hat{s}[i][j]$. Then the evaluation order induced by $\hat{s}$
is a legal full parenthesization of $A_i \cdots A_j$ for every $(i,j)$, and
$\hat{m}[1][n] \ge m[1][n]$, with equality if and only if $\hat{s}$ agrees with
some optimal split table on every sub-chain reachable from $(1,n)$.
\end{theorem}

\begin{proof}
We proceed by strong induction on the sub-chain length $\ell = j - i + 1$.

\emph{Base case} ($\ell = 1$): the sub-chain is a single matrix; the trivial
parenthesization is legal and $\hat{m}[i][i] = 0 = m[i][i]$.

\emph{Inductive step}: assume the claim holds for all sub-chains of length less
than $\ell$. Consider $(i,j)$ with $j - i + 1 = \ell$ and let
$k^\ast = \hat{s}[i][j]$. By hypothesis $i \le k^\ast < j$, so $(i, k^\ast)$
and $(k^\ast{+}1, j)$ are well-formed sub-chains of strictly smaller length.
By the inductive hypothesis, each admits a legal parenthesization induced by
$\hat{s}$; composing them with a final product of the two intermediate results
yields a legal parenthesization of $A_i \cdots A_j$ whose cost is exactly
$\hat{m}[i][j]$. Since $m[i][j]$ minimizes over all legal parenthesizations of
the sub-chain---in particular over the one induced by $\hat{s}$---we have
$\hat{m}[i][j] \ge m[i][j]$. Applying this at $(1,n)$ gives
$\hat{m}[1][n] \ge m[1][n]$.

For the equality condition: if $\hat{s}$ agrees with an optimal split table on
every sub-chain reachable from $(1,n)$ through its own splits, then unrolling
the recursion reproduces an optimal parenthesization, so
$\hat{m}[1][n] = m[1][n]$. Conversely, if $\hat{m}[1][n] = m[1][n]$, then every
split on the recursion tree must be cost-optimal for its sub-chain, since any
strictly suboptimal split at a reachable node would, by the strict additivity
of \eqref{eq:dp}, make $\hat{m}[1][n] > m[1][n]$.
\end{proof}

\begin{remark}
Theorem~\ref{thm:validity} holds for \emph{any} predictor whose output is a
split table, including an untrained one: the guarantee is architectural, not
learned. In contrast, an unconstrained regressor may output
$\hat{c} < m[1][n]$, a cost achieved by no parenthesization. We refer to the
fraction of predictions satisfying $\hat{c} \ge m[1][n]$ as
\emph{mathematical validity} and measure it empirically in Section~VI.
\end{remark}

This distinction---predicting combinatorial structure versus unconstrained
numbers---is the central axis of this study.

\section{Proposed Methodology}

\subsection{System Overview}
Fig.~1 presents the overall framework. Given a dimension sequence
$[d_0, d_1, \ldots, d_n]$, two paradigms operate in parallel. The
\emph{structural reasoning} branch deploys GraphMCMNet and PointerMCMNet to
predict split positions, from which exact costs are computed. The
\emph{statistical estimation} branch uses XGBoost and Random Forest to regress
costs from engineered features. All predictions are evaluated against
identical DP ground truth.

\subsection{GraphMCMNet: Learning the DP Structure}
The key insight behind GraphMCMNet is that the DP recurrence in \eqref{eq:dp}
naturally defines a graph over sub-problems---and a graph neural network can
learn to traverse it.

\textbf{Sub-Problem Graph Construction.} Each sub-chain $(i,j)$ for
$1 \le i \le j \le n$ becomes a node, yielding $n(n{+}1)/2$ nodes (e.g., 10
nodes for $n{=}4$, arranged in a triangular lattice from leaves to root). Three
edge types connect these nodes: (i) \emph{DP child edges} mirroring the
recurrence---connecting each parent $(i,j)$ bidirectionally to its left child
$(i,k)$ and right child $(k{+}1,j)$ for every valid split $k$; (ii)
\emph{same-length neighbor edges} linking adjacent sub-chains of identical
length, enabling lateral information flow; and (iii) implicit \emph{self-loops}
through gated residual connections. This construction gives each node direct
access to every sub-problem it depends on in the DP table.

Each node receives a compact 10-dimensional feature vector capturing its local
context: log-transformed boundary dimensions, their product, normalized
sub-chain length, mean and standard deviation of log-dimensions, min/max
statistics, the position of the minimum dimension, and a leaf indicator.

\textbf{Architecture.} GraphMCMNet (Fig.~2(b)) processes this graph in four
stages. Raw features are projected into a 128-dimensional latent space via a
linear layer with LayerNorm and SiLU activation. Next, $K{=}6$ rounds of
\emph{gated message passing} iteratively refine the embeddings: at each layer,
every node collects messages from its neighbors via a two-layer MLP,
mean-aggregates them, and updates its embedding through a learned
sigmoid-activated gate that element-wise modulates the update---preventing
over-smoothing while integrating structural information. Each layer concludes
with LayerNorm and a residual connection.

After message passing, a \emph{split scorer} decides, for each sub-chain
$(i,j)$, which split $k$ is optimal by concatenating the embeddings of the
parent node with those of the candidate children:
\begin{equation}
\mathrm{score}(k \mid i,j) = \mathrm{MLP}_{\mathrm{split}}\bigl(h_{(i,j)} \,\|\, h_{(i,k)} \,\|\, h_{(k+1,j)}\bigr)
\end{equation}
The predicted split is $\hat{s}[i][j] = \arg\max_k \mathrm{score}(k \mid i,j)$.
Scoring is fully vectorized: for each sub-chain length $L$, all sub-chains and
candidates are scored in a single batched forward pass. An auxiliary cost head
predicts $\log(1 + m[1][n])$ from the root embedding, providing a direct
gradient signal for cost accuracy during training.

\subsection{PointerMCMNet: Sequential Split Selection}
While GraphMCMNet explicitly models the DP graph, PointerMCMNet (Fig.~2(a))
processes the dimension sequence left-to-right using a Transformer encoder
\cite{vaswani2017} and ``points'' to split positions using attention
\cite{vinyals2015,bahdanau2015}. This tests whether a sequential encoder,
without explicit graph structure, can reconstruct the DP relationships
implicitly---serving as a controlled ablation of the graph inductive bias.

\textbf{Encoder.} Each of the $n{+}1$ dimension positions receives an
8-dimensional feature vector: log-dimension, products with left and right
neighbors, the triple product, normalized position, normalized magnitude, and
binary indicators for local valleys and peaks. Features are projected to 128
dimensions, enriched with sinusoidal positional encodings, and processed by a
6-layer Transformer encoder (8 heads, feed-forward dimension 512, GELU,
dropout 0.1).

\textbf{Pointer Decoder.} For each sub-chain $(i,j)$ of length $L \ge 3$, a
query concatenates the boundary hidden states $h_{i-1}$ and $h_j$ with a
32-dimensional learnable length embedding, projected through an MLP. Each
candidate split $k$ is scored with Bahdanau (additive) attention:
$\mathrm{score}(k \mid i,j) = v^{T}\tanh(W_1 q_{(i,j)} + W_2 h_k)$. The
predicted split is $\hat{s}[i][j] = i + \arg\max_k
\mathrm{softmax}(\mathrm{scores})$, and an auxiliary cost head (masked mean
pooling followed by an MLP) provides the same training signal as in
GraphMCMNet.

\subsection{Statistical Baselines}
To test whether structural reasoning is necessary, we implement strong
tree-ensemble baselines---XGBoost \cite{xgboost2016} and Random Forest
\cite{breiman2001}---armed with an extensive 145-dimensional feature vector.
If exhaustive feature engineering sufficed, these models would demonstrate it.

The 145 features span four groups: \emph{global statistics} (30 features:
chain length, dimension min/max/mean/std/median, percentiles, pattern
indicators), \emph{positional features} (8 features: extrema locations,
quartile means, trend slope), \emph{cost proxies} (6 features: log-costs from
four greedy heuristics---left-to-right, right-to-left, minimum-first, and
balanced splitting), and \emph{padded sequence features} (101 features: the
full dimension sequence and pairwise products, log-transformed and
zero-padded). Note that the cost-proxy features hand the baselines the outputs
of four greedy MCM heuristics directly, so the regression models have access
to strictly more algorithmic prior knowledge than the raw dimension sequence.

Both baselines use a dual-head ensemble: a \emph{direct head} predicting
$\log(1{+}c)$ and a \emph{ratio head} predicting $c/c_{\mathrm{greedy}}$.
Final predictions blend both:
$\hat{c} = 0.3 \cdot \hat{c}_{\mathrm{direct}} + 0.7 \cdot c_{\mathrm{greedy}}
\cdot \hat{r}$, anchoring predictions to a known-feasible greedy reference.

\subsection{Training Strategy}
\textbf{Dataset.} 120{,}000 matrix chains are generated with $n \in [5, 50]$,
balanced across four dimension distributions: \emph{Uniform}
($d_i \sim U[10, 500]$), \emph{Spiky} (alternating extremes),
\emph{Bottleneck} (near-minimal inner dimensions), and \emph{Monotone}
(linearly increasing). Each sample stores the dimension sequence, DP-optimal
cost, and the complete split table. Data is split 72\%/13\%/15\% for
training, validation, and testing, with fixed random seeds for
reproducibility.

\textbf{Curriculum Learning.} Both neural models are trained with a
multi-stage curriculum (Fig.~3, Table~\ref{tab:curriculum}) that progressively
increases chain length. Curriculum learning was chosen after preliminary runs
showed unstable gradients when long chains were introduced from the start;
stage boundaries and epoch counts were selected on the validation split.

\begin{table}[t]
\caption{Curriculum Learning Schedule}
\label{tab:curriculum}
\centering
\begin{tabular}{cccccc}
\toprule
Stage & $n_{\max}$ & Epochs & LR & Batch & Mode \\
\midrule
1 & 10 & 30 & $10^{-3}$ & 128 & Cached \\
2 & 20 & 30 & $5 \times 10^{-4}$ & 96 & Cached \\
3 & 35 & 5 & $2 \times 10^{-4}$ & 48 & Low-mem \\
4 & 50 & 8 & $10^{-4}$ & 32 & Low-mem \\
\bottomrule
\end{tabular}
\end{table}

A crucial design decision: the GNN is trained \emph{only through Stage 3}
($n_{\max}{=}35$), deliberately withholding all $n{>}35$ chains. This sets up a
rigorous test of out-of-distribution generalization: can the model compose
solutions for chains 43\% longer than anything it has seen? The Pointer
Network is trained through all four stages, so its OOD comparison in
Section~VI-B is reported alongside this caveat.

\textbf{Loss Function.} Training optimizes a composite objective balancing
structural accuracy with cost fidelity:
\begin{equation}
\mathcal{L} = 0.9 \sum_{L=2}^{n} \frac{L \cdot \mathrm{CE}_L}{\sum_{L'} L' |S_{L'}|}
+ 0.1 \cdot \mathrm{LogCosh}(\hat{y}_{\mathrm{cost}} - y_{\mathrm{cost}})
\end{equation}
The split loss uses cross-entropy weighted by sub-chain length $L$, so longer
(harder) sub-chains receive proportionally more gradient signal. The auxiliary
cost loss uses LogCosh---a smooth, outlier-robust alternative to MSE. The
$0.9/0.1$ weighting was selected on the validation split from
$\{0.5, 0.7, 0.9\}$; performance was not sensitive to this choice.

\textbf{Optimization.} AdamW ($\beta_1{=}0.9$, $\beta_2{=}0.98$, weight decay
$10^{-4}$) with cosine annealing warm restarts, gradient clipping at norm 1.0,
and mixed-precision training (AMP) on an NVIDIA RTX 4050 (6\,GB). Early
stopping (patience 15) prevents overfitting. Total training: ${\sim}25$
GPU-hours. GraphMCMNet has ${\sim}3.7$M parameters; PointerMCMNet ${\sim}5.4$M.

\section{Experimental Setup}
\label{sec:setup}

\textbf{Evaluation Protocol.} A held-out set of 500 matrix chains
(seed $= 42$, $n{=}5$--$50$) is balanced across the four distributions (125
each). Ground truth is the exact DP optimum. Inference is run on CPU (Intel,
RTX 4050 system) unless otherwise noted.

\textbf{Metrics.} Seven metrics capture complementary quality dimensions:
\emph{MAPE} (mean absolute percentage error of predicted cost vs.\ DP
optimum), \emph{median and 95th-percentile (P95) error} (robust to outliers),
\emph{$R^2$ on $\log(1{+}c)$} (correlation across magnitudes),
\emph{mathematical validity} (fraction with $\hat{c} \ge m[1][n]$;
Theorem~\ref{thm:validity}), \emph{exact match} (predicted splits identical to
DP splits; neural models only), and \emph{F1-score} at a 1\% error threshold.

\textbf{Statistical Analysis.} All pairwise model comparisons use two-sided
Wilcoxon signed-rank tests on per-chain error rates, with Cohen's $d$ for
effect size. With six pairwise comparisons, a Bonferroni-corrected threshold
of $\alpha = 0.05/6 \approx 0.0083$ applies; every reported comparison remains
significant under this correction ($p < 10^{-3}$ throughout).

\section{Results and Discussion}

\subsection{Overall Performance}
To ground the numbers: given a 5-matrix chain with dimensions
$[10, 100, 5, 50, 10, 30]$, DP finds the optimum at 10{,}500 multiplications
via $((A_1(A_2A_3))((A_4A_5)))$. GraphMCMNet recovers this exact split table.
XGBoost predicts 18{,}347---a 74.7\% overshoot. On a 40-matrix spiky chain the
GNN lands within 0.09\% of DP; Random Forest overshoots by over 200\%. Across
all 500 chains the GNN reproduces the exact DP split table 93\% of the time
(Table~\ref{tab:overall}).

\begin{table}[t]
\caption{Overall Comparison on 500 Held-Out Chains}
\label{tab:overall}
\centering
\setlength{\tabcolsep}{3.5pt}
\begin{tabular}{lcccccc}
\toprule
Model & MAPE & Median & P95 & $R^2$ & Valid & Exact \\
      & (\%) & (\%)   & (\%) & (log) & (\%) & Match \\
\midrule
GraphMCMNet   & \textbf{0.029} & \textbf{0.000} & \textbf{0.064} & \textbf{0.999999} & \textbf{100.0} & \textbf{93.0\%} \\
PointerMCMNet & 0.148 & 0.000 & 0.363 & 0.999983 & 100.0 & 88.8\% \\
XGBoost       & 36.187 & 36.408 & 76.635 & 0.9067 & 22.0 & --- \\
Random Forest & 43.284 & 39.190 & 82.863 & 0.8922 & 18.2 & --- \\
\bottomrule
\end{tabular}
\end{table}

\begin{table}[t]
\caption{Error Distribution and Threshold Accuracy}
\label{tab:thresholds}
\centering
\setlength{\tabcolsep}{3.5pt}
\begin{tabular}{lccccc}
\toprule
Model & Max Err & Spearman & $\le$1\% & $\le$5\% & $\le$10\% \\
      & (\%)    & $\rho$   &          &          &           \\
\midrule
GraphMCMNet   & \textbf{3.88} & \textbf{0.999998} & \textbf{99.2\%} & \textbf{100\%} & \textbf{100\%} \\
PointerMCMNet & 10.11 & 0.999971 & 97.0\% & 99.0\% & 99.8\% \\
XGBoost       & 149.09 & 0.9588 & 3.2\% & 14.2\% & 25.8\% \\
Random Forest & 310.16 & 0.9450 & 3.6\% & 11.6\% & 18.0\% \\
\bottomrule
\end{tabular}
\end{table}

Table~\ref{tab:thresholds} reports the error distribution beyond the mean. The
GNN's \emph{median} error is exactly zero (the majority of chains are solved
optimally), its worst single prediction over 500 chains is 3.88\%, and 100\%
of predictions fall within 5\% of the DP optimum. The tree ensembles' worst
cases reach 149\% and 310\% respectively, and fewer than 4\% of their
predictions land within 1\% of the optimum.

Both neural models maintain 100\% mathematical validity: as guaranteed by
Theorem~\ref{thm:validity}, each predicts discrete split positions, so the
resulting cost always corresponds to a real parenthesization and can never dip
below the DP lower bound. The tree ensembles violate that bound in 78--82\% of
cases. The root cause is architectural: gradient-boosted trees and random
forests treat the cost as an ordinary continuous target with no notion of the
recursive feasibility that DP enforces. The GNN holds a $5.1\times$ MAPE
advantage over the Pointer Network (0.029\% vs.\ 0.148\%) because it encodes
the $O(n^2)$ sub-problem dependencies as explicit graph edges; the Pointer
Network must reconstruct those relationships implicitly through
self-attention.

These validity numbers carry practical weight. A cost model inside a compiler
or query optimizer that routinely returns infeasible values would steer
scheduling decisions toward multiplication orders that cannot exist,
corrupting downstream passes that rely on the estimate. The 100\% validity of
the structural models eliminates this failure mode entirely.

\subsection{Generalization and Robustness}

\begin{table}[t]
\caption{MAPE (\%) Breakdown by Chain Length and Distribution}
\label{tab:breakdown}
\centering
\begin{tabular}{lcccc}
\toprule
Condition & GNN & Pointer & XGBoost & RF \\
\midrule
\multicolumn{5}{l}{\emph{By chain length}} \\
$n{=}5$--$20$ & 0.027 & 0.092 & 33.09 & 45.22 \\
$n{=}21$--$40$ & 0.036 & 0.082 & 39.57 & 43.37 \\
$n{=}41$--$50^{\dagger}$ & \textbf{0.019} & 0.391 & 35.62 & 39.23 \\
\midrule
\multicolumn{5}{l}{\emph{By dimension distribution}} \\
Uniform    & 0.003 & 0.032 & 54.18 & 51.18 \\
Spiky      & 0.114 & 0.500 & 41.40 & 46.68 \\
Bottleneck & 0.000 & 0.059 & 41.97 & 62.94 \\
Monotone   & 0.000 & 0.000 & 7.19  & 12.34 \\
\bottomrule
\multicolumn{5}{l}{\footnotesize $^{\dagger}$GNN trained on $n \le 35$ only; evaluated beyond its training horizon.}
\end{tabular}
\end{table}

Table~\ref{tab:breakdown} breaks MAPE down by chain length and dimension
distribution. The GNN was never shown chains longer than 35 during training,
yet it records its lowest error (0.019\%) on the $n{=}41$--$50$
bucket---better than its in-distribution average (0.027--0.036\%). Longer
chains are composed of many shorter sub-chains the model has already mastered;
the graph edges let it assemble solutions for larger instances out of familiar
pieces, consistent with findings in neural algorithmic reasoning that
algorithm-aligned architectures extrapolate in size
\cite{velickovic2022clrs,ibarz2022generalist}. The Pointer Network, lacking
that explicit structure, degrades $4.3\times$ on the same length bucket
(0.391\% vs.\ 0.092\% at $n{\le}20$)---although we note it was trained on the
full length range, so this comparison measures architectural robustness rather
than strict OOD extrapolation. Tree models hover between 31\% and 47\% error
at every length, reinforcing that the gap is not about scale but about the
mismatch between unconstrained regression and combinatorial structure.

On distribution difficulty: Bottleneck and Monotone chains are solved
essentially perfectly by the GNN (0.000\% MAPE at the reported precision).
Spiky chains are the hardest for every model, but the GNN's worst case there
(0.114\%) is still smaller than XGBoost's best case on any distribution
(7.19\% on Monotone). Fig.~4 visualizes the MAPE spread across distributions,
and Fig.~5 shows that structural accuracy remains high beyond the training
horizon.

\subsection{Statistical Significance}

\begin{table}[t]
\caption{Wilcoxon Signed-Rank Tests on Per-Chain Errors (Two-Sided)}
\label{tab:wilcoxon}
\centering
\setlength{\tabcolsep}{4pt}
\begin{tabular}{lccc}
\toprule
Comparison & $W$ & $p$-value & Cohen's $d$ \\
\midrule
Pointer vs.\ GNN     & 302    & $1.5 \times 10^{-4}$  & 0.13 \\
Pointer vs.\ XGBoost & 3      & $1.3 \times 10^{-83}$ & $-1.34$ \\
Pointer vs.\ RF      & 11     & $1.4 \times 10^{-83}$ & $-1.09$ \\
GNN vs.\ XGBoost     & 60     & $1.8 \times 10^{-83}$ & $-1.34$ \\
GNN vs.\ RF          & 0      & $1.3 \times 10^{-83}$ & $-1.09$ \\
XGBoost vs.\ RF      & 44{,}082 & $9.7 \times 10^{-9}$ & $-0.25$ \\
\bottomrule
\end{tabular}
\end{table}

Table~\ref{tab:wilcoxon} reports the full significance analysis. All six
pairwise comparisons are significant at $p < 0.001$, and all survive
Bonferroni correction for multiple comparisons
($\alpha = 0.0083$). The structural-vs-statistical comparisons exhibit very
large effect sizes ($|d| > 1.0$), while GNN-vs-Pointer shows a small but
significant effect ($d = 0.13$): the two structural models are close to each
other and categorically separated from the regression baselines. In
head-to-head win rates (fraction of chains where one model has strictly lower
error), the GNN wins on 100\% of chains against both tree baselines and on
91\% of chains against the Pointer Network.

\subsection{Computational Cost and Ablation}
\label{sec:latency}

\begin{table}[t]
\caption{Computational Complexity and Measured Speedup}
\label{tab:complexity}
\centering
\begin{tabular}{lccc}
\toprule
Model & Inference Complexity & Params & Speedup$^{\ast}$ \\
\midrule
DP (exact)      & $O(n^3)$   & ---   & $1\times$ \\
Hu--Shing (exact) & $O(n \log n)$ & --- & not benchmarked \\
GraphMCMNet     & $O(n^2 K)$ & 3.7\,M & ${>}16\times$ \\
PointerMCMNet   & $O(n^2)$   & 5.4\,M & ${>}16\times$ \\
XGBoost / RF    & $O(n^2)^{\ddagger}$ & --- & ${\sim}12\times$ \\
\bottomrule
\multicolumn{4}{l}{\footnotesize $^{\ast}$Wall-clock vs.\ Python DP at $n{=}50$ on CPU; ${>}100\times$ under GPU batching.} \\
\multicolumn{4}{l}{\footnotesize $^{\ddagger}$145-feature extraction dominates. $K{=}6$ message-passing rounds.}
\end{tabular}
\end{table}

Table~\ref{tab:complexity} lists computational complexity and measured
speedups. We emphasize that neural inference is \emph{not} constant time: it
scales as $O(n^2 K)$ for the GNN (with a fixed number $K$ of message-passing
rounds, i.e., constant \emph{inference depth}) and $O(n^2)$ for the Pointer
Network. The practical advantage comes from (i) a lower asymptotic order than
the $O(n^3)$ DP and (ii) full vectorization: all sub-chains of a given length
are scored in one batched pass, yielding a measured ${>}16\times$ wall-clock
speedup over Python DP on CPU at $n{=}50$ and ${>}100\times$ under GPU
batching. Two caveats apply. First, the Hu--Shing algorithm achieves
$O(n \log n)$ exactly; we did not benchmark an implementation, and for
single-chain latency at large $n$ it would likely dominate all learned
approaches---the relevant comparison for our target setting is
\emph{amortized batched throughput}, where thousands of chains are evaluated
concurrently on accelerator hardware, a regime in which sequential
combinatorial algorithms parallelize poorly. Second, our DP reference is a
Python implementation; an optimized C++ DP would shrink the measured margins.
Both caveats are revisited in Section~\ref{sec:limitations}.

Tree ensembles look fast in isolation, but building the 145-feature vector
already costs $O(n^2)$, so end-to-end they are slower than the neural models
and far less accurate.

Three ablation observations stand out, derived from the controlled
comparisons in Tables~\ref{tab:overall} and \ref{tab:breakdown}:
\begin{enumerate}
    \item \emph{Graph encoder vs.\ sequential encoder} (GraphMCMNet vs.\
    PointerMCMNet, identical training targets and loss): removing the explicit
    DP-graph inductive bias raises MAPE $5.1\times$ (0.029\% $\to$ 0.148\%)
    and degrades long-chain robustness $4.3\times$.
    \item \emph{Split prediction vs.\ cost regression} (structural vs.\
    statistical branch): removing the structural output space drops validity
    from 100\% to 18--22\%; no post-processing can restore a guarantee that
    the output space does not encode.
    \item \emph{Length extrapolation}: OOD MAPE improves relative to
    in-distribution (0.019\% vs.\ 0.027--0.036\%), evidence that the GNN has
    internalized the compositional structure of the DP recurrence rather than
    memorizing length-specific patterns.
\end{enumerate}
A finer-grained component ablation (individual edge types, curriculum stages,
auxiliary loss) is left to future work.

Taken together, on this benchmark GraphMCMNet behaves as a high-fidelity
learned approximation of the DP solver: 0.029\% error, zero validity
violations, and a $16\times$ measured speedup over the Python DP reference.
The 0.029\% residual amounts to about 3 extra multiplications per 10{,}000
optimal ones. The validity gap is arguably the most consequential finding: it
shows that regression-based cost models \cite{neo2019,balsa2022} can silently
output infeasible values in the majority of cases, a risk that structural
architectures eliminate by construction. For out-of-distribution robustness,
problem-aware graph encodings \cite{bengio2021} proved decisive.

\section{Limitations and Threats to Validity}
\label{sec:limitations}

We state the boundaries of our claims explicitly.

\textbf{Synthetic data.} All training and evaluation chains are synthetic
($d \in [10, 500]$ and structured variants). Real workloads---tensor
contraction sequences in deep learning compilers, join orders in query
optimizers---have dimension statistics that are not uniformly random and may
include dimensions $d > 10^4$. Our accuracy and validity conclusions are
therefore established for the studied distributions; validating on traces
extracted from production compilers (e.g., TVM tensor programs) is the most
important item of future work.

\textbf{No Hu--Shing runtime benchmark.} We compare against the $O(n^3)$ DP
because it is the de facto standard in the systems that motivate this work,
but a rigorous latency study must include an optimized Hu--Shing
implementation and optimized C/C++ DP as baselines. Until then, our speedup
figures should be read as ``vs.\ Python DP,'' not as a claim of dominance over
the best exact algorithms. Our accuracy and validity results are unaffected by
this caveat, since both exact algorithms produce the same optimal cost.

\textbf{Speedup measurement granularity.} We report aggregate speedup factors
rather than a per-model millisecond latency table with confidence intervals;
producing such a table (across $n$, batch sizes, and hardware) is planned.

\textbf{Scale.} Evaluation covers $n \le 50$. The GNN's sub-problem graph has
$O(n^2)$ nodes and $O(n^3)$ DP-child edges, so memory becomes the binding
constraint for $n \gg 100$; hierarchical or sparsified graph constructions
would be needed at that scale.

\textbf{Single training run.} Reported metrics come from a single trained
checkpoint per model (fixed seeds; evaluation seed 42). Variance across
training seeds, per-metric confidence intervals, and convergence curves are
not reported and would strengthen the empirical claims.

\textbf{Approximation, not replacement.} The learned models approximate the DP
solver; they do not replace it in settings requiring certified optimality.
The structural guarantee bounds predictions from below by feasibility, not
from above by optimality: the 93\% exact-match rate means 7\% of chains
receive a (slightly) suboptimal parenthesization.

\section{Conclusion}
This paper studied the Matrix Chain Multiplication problem through the lens of
learned approximation, comparing two framings under identical conditions:
structural reasoning, which predicts the DP split table and computes costs
exactly from the predicted structure, and statistical estimation, which
regresses costs directly. Evaluation across 500 diverse synthetic chains
reveals a stark dichotomy: regression-based methods fail systematically,
whereas structural neural models closely track the DP solver.

The most consequential observation is the \emph{mathematical validity gap}.
Despite 145 engineered features, tree-based regressors produce costs below the
theoretical DP minimum in roughly 80\% of cases, because they treat a
combinatorial objective as an unconstrained continuous variable. GraphMCMNet
instead predicts discrete split decisions, and---as we prove in
Theorem~\ref{thm:validity}---every such prediction corresponds to a feasible
parenthesization, yielding 100\% validity by construction. On our benchmark
this framing simultaneously reduces error to 0.029\% MAPE while providing a
measured $16\times$ speedup over the Python DP reference, with inference
scaling as $O(n^2 K)$ rather than $O(n^3)$.

Encoding DP sub-problem dependencies as graph edges also enables length
extrapolation: trained only on $n \le 35$, GraphMCMNet achieves its lowest
error on chains of length 41--50. These findings support a broader thesis from
neural algorithmic reasoning: embedding algorithmic structure directly into
neural architectures is a prerequisite for reliable learned combinatorial
optimization. Future work will evaluate on real compiler traces, benchmark
against optimized exact algorithms including Hu--Shing, report per-seed
variance and fine-grained latency, and scale the structural approach to
larger tensor computation graphs.

\begin{thebibliography}{00}
\bibitem{godbole1973} S. S. Godbole, ``On efficient computation of matrix chain products,'' \emph{IEEE Trans. Computers}, vol. C-22, no. 9, pp. 864--866, Sep. 1973.
\bibitem{hushing1982} T. C. Hu and M. T. Shing, ``Computation of matrix chain products, Part I, Part II,'' \emph{SIAM J. Computing}, vol. 11, no. 2, pp. 362--373, 1982.
\bibitem{lopez2023} F. L\'opez, L. Karlsson, and P. Bientinesi, ``On the parenthesisations of matrix chains: All are useful, few are essential,'' arXiv:2303.17352, 2023.
\bibitem{cormen2009} T. H. Cormen, C. E. Leiserson, R. L. Rivest, and C. Stein, \emph{Introduction to Algorithms}, 3rd ed. Cambridge, MA: MIT Press, 2009.
\bibitem{naz2021} A. Naz \emph{et al.}, ``Optimal sequence for chain matrix multiplication using evolutionary algorithm,'' \emph{PLOS ONE}, 2021.
\bibitem{vinyals2015} O. Vinyals, M. Fortunato, and N. Jaitly, ``Pointer networks,'' in \emph{Proc. NeurIPS}, vol. 28, 2015.
\bibitem{bello2017} I. Bello, H. Pham, Q. V. Le, M. Norouzi, and S. Bengio, ``Neural combinatorial optimization with reinforcement learning,'' in \emph{Proc. ICLR}, 2017.
\bibitem{bahdanau2015} D. Bahdanau, K. Cho, and Y. Bengio, ``Neural machine translation by jointly learning to align and translate,'' in \emph{Proc. ICLR}, 2015.
\bibitem{vaswani2017} A. Vaswani \emph{et al.}, ``Attention is all you need,'' in \emph{Proc. NeurIPS}, 2017.
\bibitem{xgboost2016} T. Chen and C. Guestrin, ``XGBoost: A scalable tree boosting system,'' in \emph{Proc. ACM SIGKDD}, pp. 785--794, 2016.
\bibitem{breiman2001} L. Breiman, ``Random forests,'' \emph{Machine Learning}, vol. 45, pp. 5--32, 2001.
\bibitem{neo2019} R. Marcus \emph{et al.}, ``Neo: A learned query optimizer,'' \emph{Proc. VLDB Endowment}, vol. 12, no. 11, pp. 1705--1718, 2019.
\bibitem{balsa2022} Z. Yang \emph{et al.}, ``Balsa: Learning a query optimizer without expert demonstrations,'' in \emph{Proc. ACM SIGMOD}, 2022.
\bibitem{cappart2021} Q. Cappart \emph{et al.}, ``Combinatorial optimization and reasoning with graph neural networks,'' in \emph{Proc. IJCAI Survey Track}, 2021.
\bibitem{bengio2021} Y. Bengio, A. Lodi, and A. Prouvost, ``Machine learning for combinatorial optimization: A methodological tour d'horizon,'' \emph{European J. Operational Research}, vol. 290, no. 2, pp. 405--421, 2021.
\bibitem{vesselinova2020} N. Vesselinova, R. Steinert, D. F. Perez-Ramirez, and M. Boman, ``Learning combinatorial optimization on graphs,'' \emph{IEEE Access}, vol. 8, pp. 120388--120416, 2020.
\bibitem{tvm2018} T. Chen \emph{et al.}, ``TVM: An automated end-to-end optimizing compiler for deep learning,'' in \emph{Proc. 13th USENIX OSDI}, 2018.
\bibitem{tiramisu2021} R. Baghdadi \emph{et al.}, ``A deep learning based cost model for automatic code optimization,'' in \emph{Proc. MLSys}, 2021.
\bibitem{velickovic2021nar} P. Veli\v{c}kovi\'c and C. Blundell, ``Neural algorithmic reasoning,'' \emph{Patterns}, vol. 2, no. 7, 2021.
\bibitem{velickovic2022clrs} P. Veli\v{c}kovi\'c \emph{et al.}, ``The CLRS algorithmic reasoning benchmark,'' in \emph{Proc. ICML}, 2022.
\bibitem{ibarz2022generalist} B. Ibarz \emph{et al.}, ``A generalist neural algorithmic learner,'' in \emph{Proc. Learning on Graphs (LoG)}, 2022.
\bibitem{bounsi2024transformers} W. Bounsi \emph{et al.}, ``Transformers meet neural algorithmic reasoners,'' arXiv:2406.09308, 2024.
\bibitem{rampasek2022gps} L. Ramp\'a\v{s}ek \emph{et al.}, ``Recipe for a general, powerful, scalable graph transformer,'' in \emph{Proc. NeurIPS}, 2022.
\bibitem{muller2023gt} L. M\"uller, M. Galkin, C. Morris, and L. Ramp\'a\v{s}ek, ``Attending to graph transformers,'' \emph{Trans. Machine Learning Research}, 2024.
\end{thebibliography}

\end{document}
